\chapter*{Introduzione}\label{ch:introduzione}
\markboth{Introduzione}{Introduzione}
\addcontentsline{toc}{section}{Introduzione}

Negli ultimi anni si sono sensibilmente espansi il dibattito e la ricerca attorno la tecnologia del \acrfull{tc}.
Con la digitalizzazione della maggior parte dei momenti della nostra vita, diventa essenziale prevedere meccanismi e garanzie di sicurezza, confidenzialità e integrità anche in ambienti informatici.
Trusted Computing -- "calcolo fidato -- si riferisce ad un tipo di tecnologia che ha l’obiettivo di garantire che dei dispositivi (come computer o telefoni cellulari) siano protetti da eventuali manomissioni come l’iniezione di codici estranei, mediante l’uso di opportuni hardware e software.
Data l’incredibile varietà del software e degli algoritmi, delle applicazioni e delle tecnologie hardware, è fondamentale che esista un metodo di riferimento per garantire affidabilità in questo settore. Ricordando è doveroso rimanere sempre consapevoli che al giorno d’oggi non esiste un meccanismo per trasformare uno standard dalla sua specifica all’implementazione in maniera impeccabile: anzi il processo di implementazione è il momento più rischioso e delicato. Un’implementazione che sarebbe altresì perfetta -- di un crittosistema perfetto, crollerebbe se esistesse anche solo un piccolo dettaglio implementativo che fornisca \enquote{backdoor} o \enquote{scorciatoie} di qualsiasi natura.  
Questo testo cerca di navigare una questione molto ampia delineando prima il problema di sicurezza in analisi, e poi una soluzione particolare ad una piattaforma. Nel dettaglio si tratterà di \acrfull{tpm} -- da qui in poi \acrshort{tpm} -- la specifica e proposte di implementazione “ufficiali”. 
Il \acrshort{tpm} è una soluzione che viene ideata in primissimo luogo pensando ai calcolatori personali, ai \acrfull{pc}, ma verranno citate anche proposte che spaziano su tutt’altro tipo di hardware, trasformandosi talvolta in soluzioni diverse allo stesso problema.
